\section{Limitations}
\label{sec:limits}

This section states what the system does not do and what we do not know. It is
placed before the roadmap because the two requests in \Cref{sec:next} follow
from it.

\subsection{State re-identification}

This is the largest known limitation in the system.

The state-identity cascade recognises a screen it has \emph{not} seen before far
more reliably than it recognises one it \emph{has}. Matching a new observation
against a state already in the graph is unreliable when the screen's contents
have changed between visits, so a revisited screen can be stored as a second
state.

Two consequences follow. Interface-state counts are inflated, so the 97 states
recorded on AUT-1 are an upper bound rather than a measurement. More
seriously, navigation memory can only shorten a route it recognises, so every
failure to re-identify is a route the agent does not reuse. This is the most
plausible explanation for why autonomy on AUT-1 recovered to 68\,\%
rather than to the 100\,\% reached on Contacts.

The fix is a content-independent structural signature rather than a comparison
over the whole accessibility tree, and it is the first piece of engineering we
would do next.

\subsection{Quality degrades sharply without a specification}

The system runs without a specification. It does not run well without one. On
AcademyBug, with nothing ingested, 16 of 20 runs were lost to the agent, against
none of AUT-2's 64. With no requirements to target, the planner has no
criterion for what the application ought to have done and no stopping condition
for a test.

The requirement as written, that no knowledge source is a hard dependency, is
satisfied. The requirement as intended is not. The remedy is a provisional
specification generated from a first exploratory pass, which has not been built.

\subsection{Deduplication is imperfect}

Findings are deduplicated against titles and a similarity threshold. Three
descriptions of one fault, generated by different runs, can fall outside that
threshold: the AUT-2 campaign reported three separate major findings for a
single control that failed to open. The effect on a reader is over-counting
rather than a false report, which is why it was ranked below the attribution
work. The fix is to deduplicate against the target element and screen state
rather than against the text.

\subsection{Measurement the project cannot supply}

Precision and recall are not reported. Every application-level failure is a
candidate defect requiring human confirmation, and the system is scoped to say
so. Without a build whose defects are known in advance there is no way to
compute either figure, and the broken-workshop result, zero defects found on a
deliberately faulty application, is the clearest illustration of the gap: we
cannot tell whether the faults were unreachable or simply not provoked.

\subsection{Configuration was not held fixed}

The model assignment changed eight times across seven vendors between April and
September. Where a campaign changed both an agent mechanism and a model, the
credit between them is not cleanly separable. Pass rates rose over the period
and so did the quality of available models, and this report does not separate
those two causes.

\subsection{Scale not yet demonstrated}

The largest graph built is roughly three orders of magnitude below the
ETA-NFR-002 ceilings. No degradation was observed in the range exercised, which
is not the same as demonstrating the ceiling. Campaigns were held at 25 tests
deliberately, because that is what one engineer can review by hand, and because
the planner's duplicate-suppression memory saturates near that scale with the
current caps.

\subsection{Positioning against published work}
\label{sec:positioning}

Seven comparable language-model testing systems were reviewed in full. Of those,
five are handed their test objective and only one generates its own. None
distinguishes an agent failure from an application defect from an environment
fault, and none reports precision or recall for defect detection. The two works
in the wider literature that do report those figures injected their faults by
hand.

The gap is therefore the field's and not only this project's, which is context
for the request that follows rather than an excuse. A system grounded in a
specification attacks the oracle problem directly, and that is precisely what
makes a seeded-defect experiment on this system worth running: it is one of the
few setups where the answer would mean something.
